\begin{explanationbox}
	\textbf{\textcolor{CExplanation}{一句话直觉}}

	正四面体的高和棱长的比值是确定的常数，类似于等边三角形高与边长的关系。

	\medskip
	\textbf{\textcolor{CExplanation}{核心拆解}}
	\begin{description}[style=nextline, leftmargin=2.8em, labelsep=0.5em, itemsep=0.45em]
		\item[\textbf{要点一（底面中心定位）：}]正四面体的顶点在底面的投影是底面正三角形的中心（重心、外心合一）。
		\item[\textbf{要点二（直角三角形勾股）：}]高、侧棱、底面中心到顶点距离构成直角三角形。
		\item[\textbf{要点三（比例定值）：}]计算得此直角三角形中，侧棱为斜边，另一直角边为$a/\sqrt{3}$，高为$\sqrt{6}/3 a$.
	\end{description}

	\medskip
	\textbf{\textcolor{CExplanation}{几何本质（空间对称性）}}

	正四面体是空间中最对称的多面体之一，顶点在底面的投影恰好是底面正三角形的中心，这使得高线成为对称轴的一部分。

	\medskip
	\textbf{\textcolor{CExplanation}{代数意义（定量计算）}}

	该公式将几何度量转化为代数计算，一旦已知棱长，可立即得到高，进而计算体积、表面积、内切球半径等。

	\medskip
	\textbf{\textcolor{CExplanation}{考点价值（高频应用）}}
	\begin{itemize}[leftmargin=*, itemsep=0.5em, topsep=0.4em]
		\item \textbf{考法一（直接套用）：}已知棱长求高，直接代入公式。
		\item \textbf{考法二（反求棱长）：}已知高求棱长或其他量，先通过公式解出棱长。
	\end{itemize}

	\medskip
	\textbf{\textcolor{CExplanation}{顿悟点}}

	将空间立体问题转化为平面直角三角形问题，是解决正四面体问题的核心思想。

	\medskip
	\textbf{\textcolor{CExplanation}{使用场景}}
	\begin{itemize}[leftmargin=*, itemsep=0.5em, topsep=0.4em]
		\item \textbf{场景一（基础计算）：}已知棱长求高或已知高求棱长。
		\item \textbf{场景二（综合计算）：}在体积、表面积、球半径等问题中作为中间量。
	\end{itemize}
\end{explanationbox}
